% =============================================
% RESUMO
% =============================================
\newpage
\thispagestyle{empty}
\begin{center}
    \fontsize{14pt}{14pt}\selectfont\bfseries RESUMO\\[24pt]
\end{center}

\noindent LONARDON JÚNIOR, M. A.; BONADIO, M. L.; MOREIRA, V. N. Título. 2026. Número de páginas. Trabalho de Conclusão de Curso (Bacharelado em Engenharia de Software) -- Centro de Ciências Exatas, Naturais e Tecnológicas, Universidade de Ribeirão Preto, Ribeirão Preto -- SP, 2026.\\[12pt]

A gestão financeira pessoal no Brasil ainda apresenta limitações quanto ao acesso a ferramentas capazes de transformar dados bancários em informações úteis para a tomada de decisão cotidiana. Embora o Banco Central do Brasil disponibilize o \textit{Open Finance}, uma ferramenta que possibilita o compartilhamento padronizado de dados financeiros entre instituições autorizadas, grande parte das soluções disponíveis concentra-se na visualização estática de informações, com baixa capacidade de interpretação contextual e antecipação de cenários. Diante disso, este trabalho propõe o desenvolvimento de uma solução de assistência financeira inteligente operada integralmente por meio do WhatsApp, sem dependência de plataforma \textit{web}, com foco em acessibilidade, praticidade e interação conversacional. A proposta integra dados financeiros obtidos via \textit{Open Finance} a técnicas de análise preditiva, aprendizado de máquina e séries temporais, com o objetivo de identificar padrões de comportamento, gerar projeções personalizadas e emitir alertas sobre possíveis riscos ou oportunidades financeiras. Complementarmente, utiliza-se um assistente conversacional baseado em \textit{Retrieval-Augmented Generation} (RAG), permitindo ao usuário consultar, em linguagem natural, seu histórico financeiro, suas projeções e recomendações associadas ao seu contexto individual. Como resultado esperado, pretende-se oferecer uma alternativa mais fluida e próxima da rotina do usuário, centralizando a experiência em um canal amplamente utilizado e reduzindo barreiras de adoção tecnológica. Espera-se, assim, ampliar a compreensão sobre a própria vida financeira, favorecer o planejamento pessoal e apoiar decisões mais conscientes por meio de uma interação simples, contínua e inteligente.\\

\noindent\textbf{Palavras-chave:} \textit{Open Finance}; WhatsApp; Inteligência artificial; Análise preditiva; Assistência financeira conversacional.

% =============================================
% ABSTRACT
% =============================================
\newpage
\thispagestyle{empty}
\begin{center}
    \fontsize{14pt}{14pt}\selectfont\bfseries ABSTRACT\\[24pt]
\end{center}

\noindent LONARDON JÚNIOR, M. A.; BONADIO, M. L.; MOREIRA, V. N. Title in English. 2026. Number of pages (pages). Undergraduate Thesis (Bachelor's degree of Software Engineering) -- Center of Exact, Natural and Technological Sciences, University of Ribeirão Preto, Ribeirão Preto -- SP, 2026.\\[12pt]

Personal financial management in Brazil still faces limitations regarding access to tools capable of transforming banking data into useful information for everyday decision-making. Although the Central Bank of Brazil provides \textit{Open Finance}, a framework that enables the standardized sharing of financial data among authorized institutions, most available solutions focus on the static visualization of information, with limited capacity for contextual interpretation and scenario forecasting. In light of this, this study proposes the development of an intelligent financial assistance solution operated entirely through WhatsApp, without reliance on a web platform, focusing on accessibility, practicality, and conversational interaction. The proposed solution integrates financial data obtained through \textit{Open Finance} with predictive analysis techniques, machine learning, and time series analysis in order to identify behavioral patterns, generate personalized forecasts, and issue alerts regarding potential financial risks or opportunities. Additionally, a conversational assistant based on \textit{Retrieval-Augmented Generation} (RAG) is employed, allowing users to query their financial history, forecasts, and recommendations associated with their individual context using natural language. As an expected outcome, the project aims to provide a more fluid and user-centered alternative by centralizing the experience within a widely used communication channel and reducing barriers to technological adoption. Thus, the study seeks to enhance users' understanding of their own financial lives, support personal financial planning, and encourage more informed decision-making through a simple, continuous, and intelligent interaction.\\[12pt]

\noindent\textbf{Keywords:} \textit{Open Finance}; WhatsApp; Artificial intelligence; Predictive analytics; Conversational financial assistance.
